\section{Máquina de estados da Unidade de Controle}
\label{sec:fsm}

A UC é uma máquina de Moore com quatro estados, percorridos sempre na mesma
ordem por todas as instruções (Figura~\ref{fig:fsm}). O estado só muda nas
bordas de subida em que \sig{CE}~=~1. O reset vem da chave SW0 (sincronizada por
dois flip-flops) e, de qualquer estado, leva a máquina a FETCH com $PC = 0$,
$SP = 254$, flags e registradores zerados; ao levar a chave para 0, o programa
começa a ser executado a partir do endereço 0.

\begin{figure}[H]
\centering
\begin{tikzpicture}[
  st/.style={state, thick, minimum size=2.3cm, align=center, font=\small\bfseries, fill=blue!6},
  ->, >={Latex[length=2.5mm]}, thick, node distance=4.2cm]

  \node[st, initial, initial text={\scriptsize reset}] (f) {FETCH};
  \node[st, right=of f] (d) {DECODE};
  \node[st, below=2.6cm of d] (e) {EXECUTE};
  \node[st, below=2.6cm of f] (w) {WRITE};

  \draw (f) -- node[above, font=\scriptsize] {CE} (d);
  \draw (d) -- node[right, font=\scriptsize] {CE} (e);
  \draw (e) -- node[below, font=\scriptsize] {CE} (w);
  \draw (w) -- node[left, font=\scriptsize] {CE} (f);

  \node[font=\scriptsize, align=left, anchor=south, text width=4.6cm] at ([yshift=2pt]f.north)
       {IR $\gets$ DOUT (= MEM[PC])\\MAR $\gets$ PC+1};
  \node[font=\scriptsize, align=left, anchor=south, text width=4.6cm] at ([yshift=2pt]d.north)
       {decodifica IR; ALU\_A, ALU\_B, CMD\\MBR $\gets$ 2º byte (\instr{ldi}, \instr{jmp})\\MAR/MBR/WE para acesso à memória};
  \node[font=\scriptsize, align=left, anchor=north, text width=4.6cm] at ([yshift=-2pt]e.south)
       {MBR $\gets$ ALU\_R e FLAGS $\gets$ flags (ALU)\\MBR $\gets$ DOUT (\instr{ldr}, \instr{pop}); MBR $\gets$ Ry (\instr{mov})\\escrita na RAM; WE $\gets$ 0};
  \node[font=\scriptsize, align=left, anchor=north, text width=4.6cm] at ([yshift=-2pt]w.south)
       {Rx $\gets$ MBR; SP $\pm$ 1\\PC $\gets$ próximo PC\\MAR $\gets$ próximo PC};
\end{tikzpicture}
\caption{Máquina de estados da Unidade de Controle.}
\label{fig:fsm}
\end{figure}

\subsection{Descrição dos estados}

\paragraph{FETCH.} Como o \sig{MAR} já contém o \sig{PC} (carregado no WRITE
anterior ou no reset), \sig{DOUT} já mostra $MEM[PC]$ e a instrução é copiada
para o \sig{IR}. Ao mesmo tempo, $MAR \gets PC + 1$, antecipando a leitura do
segundo byte. Para instruções de um byte essa leitura é simplesmente ignorada;
fazê-la sempre simplifica a UC e mantém o número de estados constante.

\paragraph{DECODE.} Com o \sig{IR} válido, a instrução é decodificada e
\sig{DOUT} mostra $MEM[PC+1]$. Neste estado:
\begin{itemize}[nosep]
  \item instruções de ALU carregam \sig{ALU\_A}, \sig{ALU\_B} e \sig{CMD};
  \item \instr{ldi} e \instr{jmp} copiam o segundo byte para \sig{MBR};
  \item \instr{st} usa o segundo byte como endereço ($MAR \gets DOUT$),
        $MBR \gets Rx$ e $WE \gets 1$;
  \item \instr{str} e \instr{push} preparam endereço e dado e ligam \sig{WE};
  \item \instr{ldr} e \instr{pop} colocam o endereço de leitura em \sig{MAR};
  \item os registradores de exibição do LCD são atualizados com o \sig{IR} e o
        segundo byte (Seção~\ref{sec:lcd}).
\end{itemize}

\paragraph{EXECUTE.} A ALU já calculou o resultado (ela é combinacional e suas
entradas estão estáveis desde o fim do DECODE) e a RAM já respondeu à leitura
pedida. Nas instruções de ALU, o resultado vai para \sig{MBR} e as flags para os
flip-flops \sig{FLAGS}; em \instr{ldr} e \instr{pop}, o dado lido vai para
\sig{MBR}; em \instr{mov}, \sig{Ry} vai para \sig{MBR}. As escritas na memória
acontecem durante este estado, e \sig{WE} é desligado no fim dele.

\paragraph{WRITE.} O resultado é escrito no registrador destino
($Rx \gets MBR$), o \sig{SP} é atualizado (\instr{push}/\instr{pop}), e o
próximo \sig{PC} é calculado e escrito tanto em \sig{PC} quanto em \sig{MAR},
deixando a memória pronta para o próximo FETCH. As condições dos saltos são
avaliadas aqui, a partir das flags salvas --- que, nos saltos, nunca mudam
durante a própria instrução.

\subsection{Tempo de execução}

Com quatro estados de \SI{0,5}{\second}, toda instrução leva exatamente
\SI{2}{\second}, inclusive as de dois bytes, os saltos e \instr{halt}. O tempo
constante facilita o acompanhamento visual e a montagem da tabela passo a
passo. No \instr{halt}, o próximo PC é o próprio PC: a UC continua percorrendo
os quatro estados, mas sempre reexecuta \instr{halt}, sem alterar nada, até que
o reset seja acionado.

\subsection{Micro-operações por instrução}

A Tabela~\ref{tab:microops} resume o que cada classe de instrução faz em cada
estado. O FETCH é igual para todas ($IR \gets DOUT$, $MAR \gets PC+1$) e foi
omitido.

\begin{table}[H]
\centering
\caption{Micro-operações por instrução nos estados DECODE, EXECUTE e WRITE.}
\label{tab:microops}
\footnotesize
\renewcommand{\arraystretch}{1.25}
\begin{tabularx}{\linewidth}{>{\raggedright}p{2.3cm} >{\raggedright}X >{\raggedright}X >{\raggedright\arraybackslash}X}
\toprule
Instrução & DECODE & EXECUTE & WRITE \\
\midrule
\instr{add}, \instr{sub}, \instr{and}, \instr{or}, \instr{xor}
  & $A \gets Rx$; $B \gets Ry$; CMD
  & $MBR \gets ALU\_R$; $FLAGS \gets$ flags
  & $Rx \gets MBR$; $PC \gets PC+1$ \\
\instr{inc}, \instr{dec}
  & $A \gets Rx$; $B \gets 1$; CMD
  & $MBR \gets ALU\_R$; $FLAGS \gets$ flags
  & $Rx \gets MBR$; $PC \gets PC+1$ \\
\instr{not}, \instr{rol}, \instr{ror}, \instr{lsl}, \instr{lsr}
  & $A \gets Rx$; $B \gets 0$; CMD
  & $MBR \gets ALU\_R$; $FLAGS \gets$ flags
  & $Rx \gets MBR$; $PC \gets PC+1$ \\
\midrule
\instr{ldi Rx, v}
  & $MBR \gets DOUT$ ($=v$)
  & ---
  & $Rx \gets MBR$; $PC \gets PC+2$ \\
\instr{st Rx, a}
  & $MAR \gets DOUT$ ($=a$); $MBR \gets Rx$; $WE \gets 1$
  & escrita; $WE \gets 0$
  & $PC \gets PC+2$ \\
\instr{ldr Rx, [Ry]}
  & $MAR \gets Ry$
  & $MBR \gets DOUT$
  & $Rx \gets MBR$; $PC \gets PC+1$ \\
\instr{str Rx, [Ry]}
  & $MAR \gets Ry$; $MBR \gets Rx$; $WE \gets 1$
  & escrita; $WE \gets 0$
  & $PC \gets PC+1$ \\
\instr{push Rx}
  & $MAR \gets SP$; $MBR \gets Rx$; $WE \gets 1$
  & escrita; $WE \gets 0$
  & $SP \gets SP-1$; $PC \gets PC+1$ \\
\instr{pop Rx}
  & $MAR \gets SP+1$
  & $MBR \gets DOUT$
  & $Rx \gets MBR$; $SP \gets SP+1$; $PC \gets PC+1$ \\
\instr{mov Rx, Ry}
  & ---
  & $MBR \gets Ry$
  & $Rx \gets MBR$; $PC \gets PC+1$ \\
\midrule
\instr{jmp a}
  & $MBR \gets DOUT$ ($=a$)
  & ---
  & $PC \gets MBR$ \\
\instr{jmpr Rx}
  & ---
  & ---
  & $PC \gets Rx$ \\
\instr{bz}, \instr{bnz}, \ldots, \instr{bneg}
  & ---
  & ---
  & se condição: $PC \gets Rx$, senão $PC \gets PC+1$ \\
\midrule
\instr{nop} (e inválidas)
  & ---
  & ---
  & $PC \gets PC+1$ \\
\instr{halt}
  & ---
  & ---
  & $PC \gets PC$ \\
\bottomrule
\end{tabularx}
\end{table}

Em todas as instruções, o WRITE também faz $MAR \gets$ próximo PC, e o DECODE
atualiza os registradores de exibição do LCD. A Tabela~\ref{tab:cond} lista as
condições dos saltos.

\begin{table}[H]
\centering
\caption{Condições dos saltos (flags guardadas na UC pela última instrução de ALU).}
\label{tab:cond}
\small
\begin{tabular}{l l @{\hspace{1.5cm}} l l}
\toprule
Instrução & Salta se & Instrução & Salta se \\
\midrule
\instr{bz}   & ZERO = 1      & \instr{bneq} & EQUAL = 0 \\
\instr{bnz}  & ZERO = 0      & \instr{bgt}  & GREATER = 1 \\
\instr{bcs}  & OVERFLOW = 1  & \instr{blt}  & SMALLER = 1 \\
\instr{bcc}  & OVERFLOW = 0  & \instr{bneg} & NEGATIVE = 1 \\
\instr{beq}  & EQUAL = 1     & \instr{jmpr} & sempre \\
\bottomrule
\end{tabular}
\end{table}
